\documentclass[12pt]{article}
\usepackage[utf8]{inputenc}
\usepackage[T1]{fontenc}
\usepackage{hyperref}
\title{Yorta Yorta Language Resources for K-12 Curriculum Development}
\author{Compiled by Claude Sonnet 5.5 (claude-sonnet-5-5) for Andr\'as Kornai}
\date{}
\begin{document}
\maketitle

\section{Introduction}
Yorta Yorta (ISO 639-3 xyy; Pama-Nyungan, Victoria, Australia) has no fluent first-language speakers today. It is being revived by the Yorta Yorta community. This report covers only what appeared in web search results. Little of it is detailed, so it should be checked with the community organisations named below.

\section{Primary Educational Resources}
\begin{itemize}
\item Gowrie Street Primary School, Shepparton, began teaching Yorta Yorta in 2015. A community-led program reportedly reaches students in Years P--6. This comes from a secondary description: \url{https://vaeai.org.au/?p=6888}
\item Kaiela Institute case study. It describes the Yorta Yorta language revival work, including the Yalca Loitchbja (Talk 'N Tea) program: \url{https://www.jawun.org.au/s/GM_Case-Study_Kaiela-Institute_Yorta-Yorta-Language.pdf}
\end{itemize}

\section{Linguistic and Reference Materials}
\begin{itemize}
\item Overview article: \url{https://en.wikipedia.org/wiki/Yorta_Yorta_language}
\item ``Yorta Yorta Language Heritage'' (Peeler, Atkinson and Bowe, Monash University) is reported to include introductory lessons and English--Yorta Yorta and Yorta Yorta--English dictionaries. Search results gave no direct link. The report appears to be described at \url{https://sheppartoninterfaith.org.au/?p=2560}
\end{itemize}

\section{Religious and Community Resources}
\begin{itemize}
\item Kaiela Institute (see the case study above) is the community organisation behind the revival program.
\end{itemize}

\section{Digital and Audio Resources}
\begin{itemize}
\item ``Yorta Yorta'' dictionary app. It is reported to offer searchable word lists in both directions, audio pronunciations and offline use. Listings disagree on its size, from over 500 to over 1500 entries: \url{https://apps.apple.com/app/id1356419603}
\item App flyer: \url{https://www.moira.vic.gov.au/files/sharedassets/public/v/1/01-residents/maternal-and-child-health/yorta-yorta-app-flyer.docx}
\end{itemize}

\section{Curriculum Development Context}
School teaching exists at one primary school and appears to be community-led. No graded textbook or curriculum documents were found in the search results. A K-12 curriculum would need to build on the dictionary app and the Monash materials, and to be developed with Yorta Yorta Elders and language workers.
\end{document}